All results are on the test sessions, which were not used for training, model selection, or calibration. ``Initial'' is the initial Seq2Seq-GRU retrained on the corrected data with the same protocol as the final model. Unless stated otherwise, numbers are means over three seeds and changes are relative to the Initial seed mean. All \NRuns{} neural runs stopped early after \EpochsMin{} to \EpochsMax{} epochs (median \EpochsMedian{}). The selected checkpoints therefore lie at or before epoch \BestEpochMax{}, well inside the 60-epoch budget.

\subsection{Headline: RTT}
Table~\ref{tab:main} gives the RTT results per link and Figures~\ref{fig:mae}--\ref{fig:mis} show them per horizon step. Against the retrained initial model, the final model:
\begin{itemize}[leftmargin=*]
  \item lowers the RTT MAE on all six links, by \ChgMAEWorst\% to \ChgMAEBest\% (mean \ChgMAEMean\%). The largest gains are on 5G UL (\ChgMAEQUL\%), Starlink DL (\ChgMAESLDL\%), and OneWeb DL (\ChgMAEOWDL\%).
  \item lowers the calibrated 95\% interval score on all six links, by \ChgMISWorst\% to \ChgMISBest\% (mean \ChgMISMean\%), at the same coverage. Measured test coverage is \PICPMin{} to \PICPMax{} for Final and similar for Initial. Both intervals are calibrated, and the Final intervals are narrower because their width follows the input (Figure~\ref{fig:trace}).
  \item leaves RMSE and $R^2$ nearly unchanged: the RMSE changes by \ChgRMSEBest\% to $+$\ChgRMSEWorst\% (mean \ChgRMSEMean\%).
\end{itemize}
With three-seed ensembles on both sides, the MAE changes by \ChgEnsMAEWorst\% to \ChgEnsMAEBest\%, the RMSE by \ChgEnsRMSEWorst\% to \ChgEnsRMSEBest\%, and the interval score by \ChgEnsMISWorst\% to \ChgEnsMISBest\%. 
\begin{table}[t]
\centering
\caption{RTT results on the test sessions, averaged over the 15 steps (0.2--3\,s). Mean $\pm$ standard deviation over 3 seeds, and 3-seed ensembles. MAE, RMSE and MIS in ms. MIS is the 95\% interval score with intervals calibrated on the validation sessions (split conformal for Initial, CQR for Final). PICP is the measured coverage of those intervals.}
\label{tab:main}
\scriptsize
\setlength{\tabcolsep}{3pt}
\begin{tabular}{llccccc}
\toprule
Link & Model & MAE & RMSE & $R^2$ & MIS & PICP \\
\midrule
\tabinput{data/tab_main.tex}
\bottomrule
\end{tabular}
\end{table}

\begin{figure}[p]
  \centering
  \resizebox{\linewidth}{!}{\begin{tikzpicture}
\begin{groupplot}[group style={group size=3 by 2, horizontal sep=1.3cm, vertical sep=1.5cm}, frbase, width=5.3cm, height=4.0cm, xmin=0.2, xmax=3, xtick={0.2,1,2,3}, xlabel={horizon (s)}]
\nextgroupplot[title={OneWeb DL}, ylabel={RTT MAE (ms)}, legend to name=legMAE, legend columns=5]
\addplot[mRidge] table[x expr=\thisrow{h}*0.2, y=ridgefr] {data/h_oneweb_dl_rtt_MAE.dat};
\addlegendentry{Ridge-E}
\addplot[mPrev] table[x expr=\thisrow{h}*0.2, y=prevseq2seqens] {data/h_oneweb_dl_rtt_MAE.dat};
\addlegendentry{Initial Seq2Seq (3-seed ens.)}
\addplot[mResid] table[x expr=\thisrow{h}*0.2, y=frdist] {data/h_oneweb_dl_rtt_MAE.dat};
\addlegendentry{Final (seed mean)}
\addplot[mFR] table[x expr=\thisrow{h}*0.2, y=frdistens] {data/h_oneweb_dl_rtt_MAE.dat};
\addlegendentry{Final (3-seed ens.)}
\nextgroupplot[title={Starlink DL}]
\addplot[mRidge] table[x expr=\thisrow{h}*0.2, y=ridgefr] {data/h_starlink_dl_rtt_MAE.dat};
\addplot[mPrev] table[x expr=\thisrow{h}*0.2, y=prevseq2seqens] {data/h_starlink_dl_rtt_MAE.dat};
\addplot[mResid] table[x expr=\thisrow{h}*0.2, y=frdist] {data/h_starlink_dl_rtt_MAE.dat};
\addplot[mFR] table[x expr=\thisrow{h}*0.2, y=frdistens] {data/h_starlink_dl_rtt_MAE.dat};
\nextgroupplot[title={5G DL}]
\addplot[mRidge] table[x expr=\thisrow{h}*0.2, y=ridgefr] {data/h_quectel_dl_rtt_MAE.dat};
\addplot[mPrev] table[x expr=\thisrow{h}*0.2, y=prevseq2seqens] {data/h_quectel_dl_rtt_MAE.dat};
\addplot[mResid] table[x expr=\thisrow{h}*0.2, y=frdist] {data/h_quectel_dl_rtt_MAE.dat};
\addplot[mFR] table[x expr=\thisrow{h}*0.2, y=frdistens] {data/h_quectel_dl_rtt_MAE.dat};
\nextgroupplot[title={OneWeb UL}, ylabel={RTT MAE (ms)}]
\addplot[mRidge] table[x expr=\thisrow{h}*0.2, y=ridgefr] {data/h_oneweb_ul_rtt_MAE.dat};
\addplot[mPrev] table[x expr=\thisrow{h}*0.2, y=prevseq2seqens] {data/h_oneweb_ul_rtt_MAE.dat};
\addplot[mResid] table[x expr=\thisrow{h}*0.2, y=frdist] {data/h_oneweb_ul_rtt_MAE.dat};
\addplot[mFR] table[x expr=\thisrow{h}*0.2, y=frdistens] {data/h_oneweb_ul_rtt_MAE.dat};
\nextgroupplot[title={Starlink UL}]
\addplot[mRidge] table[x expr=\thisrow{h}*0.2, y=ridgefr] {data/h_starlink_ul_rtt_MAE.dat};
\addplot[mPrev] table[x expr=\thisrow{h}*0.2, y=prevseq2seqens] {data/h_starlink_ul_rtt_MAE.dat};
\addplot[mResid] table[x expr=\thisrow{h}*0.2, y=frdist] {data/h_starlink_ul_rtt_MAE.dat};
\addplot[mFR] table[x expr=\thisrow{h}*0.2, y=frdistens] {data/h_starlink_ul_rtt_MAE.dat};
\nextgroupplot[title={5G UL}]
\addplot[mRidge] table[x expr=\thisrow{h}*0.2, y=ridgefr] {data/h_quectel_ul_rtt_MAE.dat};
\addplot[mPrev] table[x expr=\thisrow{h}*0.2, y=prevseq2seqens] {data/h_quectel_ul_rtt_MAE.dat};
\addplot[mResid] table[x expr=\thisrow{h}*0.2, y=frdist] {data/h_quectel_ul_rtt_MAE.dat};
\addplot[mFR] table[x expr=\thisrow{h}*0.2, y=frdistens] {data/h_quectel_ul_rtt_MAE.dat};
\end{groupplot}
\node at ($(group c2r2.south)+(0,-1.25cm)$) {\pgfplotslegendfromname{legMAE}};
\end{tikzpicture}
}
  \caption{RTT MAE per horizon step. Final uses its median output. Lower is better.}
  \label{fig:mae}
\end{figure}

\begin{figure}[p]
  \centering
  \resizebox{\linewidth}{!}{\begin{tikzpicture}
\begin{groupplot}[group style={group size=3 by 2, horizontal sep=1.3cm, vertical sep=1.5cm}, frbase, width=5.3cm, height=4.0cm, xmin=0.2, xmax=3, xtick={0.2,1,2,3}, xlabel={horizon (s)}]
\nextgroupplot[title={OneWeb DL}, ylabel={RTT RMSE (ms)}, legend to name=legRMSE, legend columns=5]
\addplot[mRidge] table[x expr=\thisrow{h}*0.2, y=ridgefr] {data/h_oneweb_dl_rtt_RMSE.dat};
\addlegendentry{Ridge-E}
\addplot[mPrev] table[x expr=\thisrow{h}*0.2, y=prevseq2seqens] {data/h_oneweb_dl_rtt_RMSE.dat};
\addlegendentry{Initial Seq2Seq (3-seed ens.)}
\addplot[mResid] table[x expr=\thisrow{h}*0.2, y=frdist] {data/h_oneweb_dl_rtt_RMSE.dat};
\addlegendentry{Final (seed mean)}
\addplot[mFR] table[x expr=\thisrow{h}*0.2, y=frdistens] {data/h_oneweb_dl_rtt_RMSE.dat};
\addlegendentry{Final (3-seed ens.)}
\nextgroupplot[title={Starlink DL}]
\addplot[mRidge] table[x expr=\thisrow{h}*0.2, y=ridgefr] {data/h_starlink_dl_rtt_RMSE.dat};
\addplot[mPrev] table[x expr=\thisrow{h}*0.2, y=prevseq2seqens] {data/h_starlink_dl_rtt_RMSE.dat};
\addplot[mResid] table[x expr=\thisrow{h}*0.2, y=frdist] {data/h_starlink_dl_rtt_RMSE.dat};
\addplot[mFR] table[x expr=\thisrow{h}*0.2, y=frdistens] {data/h_starlink_dl_rtt_RMSE.dat};
\nextgroupplot[title={5G DL}]
\addplot[mRidge] table[x expr=\thisrow{h}*0.2, y=ridgefr] {data/h_quectel_dl_rtt_RMSE.dat};
\addplot[mPrev] table[x expr=\thisrow{h}*0.2, y=prevseq2seqens] {data/h_quectel_dl_rtt_RMSE.dat};
\addplot[mResid] table[x expr=\thisrow{h}*0.2, y=frdist] {data/h_quectel_dl_rtt_RMSE.dat};
\addplot[mFR] table[x expr=\thisrow{h}*0.2, y=frdistens] {data/h_quectel_dl_rtt_RMSE.dat};
\nextgroupplot[title={OneWeb UL}, ylabel={RTT RMSE (ms)}]
\addplot[mRidge] table[x expr=\thisrow{h}*0.2, y=ridgefr] {data/h_oneweb_ul_rtt_RMSE.dat};
\addplot[mPrev] table[x expr=\thisrow{h}*0.2, y=prevseq2seqens] {data/h_oneweb_ul_rtt_RMSE.dat};
\addplot[mResid] table[x expr=\thisrow{h}*0.2, y=frdist] {data/h_oneweb_ul_rtt_RMSE.dat};
\addplot[mFR] table[x expr=\thisrow{h}*0.2, y=frdistens] {data/h_oneweb_ul_rtt_RMSE.dat};
\nextgroupplot[title={Starlink UL}]
\addplot[mRidge] table[x expr=\thisrow{h}*0.2, y=ridgefr] {data/h_starlink_ul_rtt_RMSE.dat};
\addplot[mPrev] table[x expr=\thisrow{h}*0.2, y=prevseq2seqens] {data/h_starlink_ul_rtt_RMSE.dat};
\addplot[mResid] table[x expr=\thisrow{h}*0.2, y=frdist] {data/h_starlink_ul_rtt_RMSE.dat};
\addplot[mFR] table[x expr=\thisrow{h}*0.2, y=frdistens] {data/h_starlink_ul_rtt_RMSE.dat};
\nextgroupplot[title={5G UL}]
\addplot[mRidge] table[x expr=\thisrow{h}*0.2, y=ridgefr] {data/h_quectel_ul_rtt_RMSE.dat};
\addplot[mPrev] table[x expr=\thisrow{h}*0.2, y=prevseq2seqens] {data/h_quectel_ul_rtt_RMSE.dat};
\addplot[mResid] table[x expr=\thisrow{h}*0.2, y=frdist] {data/h_quectel_ul_rtt_RMSE.dat};
\addplot[mFR] table[x expr=\thisrow{h}*0.2, y=frdistens] {data/h_quectel_ul_rtt_RMSE.dat};
\end{groupplot}
\node at ($(group c2r2.south)+(0,-1.25cm)$) {\pgfplotslegendfromname{legRMSE}};
\end{tikzpicture}
}
  \caption{RTT RMSE per horizon step. Final uses its mean output.}
  \label{fig:rmse}
\end{figure}

\begin{figure}[t]
  \centering
  \resizebox{\linewidth}{!}{\begin{tikzpicture}
\begin{groupplot}[group style={group size=3 by 2, horizontal sep=1.3cm, vertical sep=1.5cm}, frbase, width=5.3cm, height=4.0cm, xmin=0.2, xmax=3, xtick={0.2,1,2,3}, xlabel={horizon (s)}]
\nextgroupplot[title={OneWeb DL}, ylabel={RTT interval score (ms)}, legend to name=legMIS, legend columns=5]
\addplot[mRidge] table[x expr=\thisrow{h}*0.2, y=ridgefr] {data/h_oneweb_dl_rtt_MIS.dat};
\addlegendentry{Ridge-E}
\addplot[mPrev] table[x expr=\thisrow{h}*0.2, y=prevseq2seqens] {data/h_oneweb_dl_rtt_MIS.dat};
\addlegendentry{Initial Seq2Seq (3-seed ens.)}
\addplot[mResid] table[x expr=\thisrow{h}*0.2, y=frdist] {data/h_oneweb_dl_rtt_MIS.dat};
\addlegendentry{Final (seed mean)}
\addplot[mFR] table[x expr=\thisrow{h}*0.2, y=frdistens] {data/h_oneweb_dl_rtt_MIS.dat};
\addlegendentry{Final (3-seed ens.)}
\nextgroupplot[title={Starlink DL}]
\addplot[mRidge] table[x expr=\thisrow{h}*0.2, y=ridgefr] {data/h_starlink_dl_rtt_MIS.dat};
\addplot[mPrev] table[x expr=\thisrow{h}*0.2, y=prevseq2seqens] {data/h_starlink_dl_rtt_MIS.dat};
\addplot[mResid] table[x expr=\thisrow{h}*0.2, y=frdist] {data/h_starlink_dl_rtt_MIS.dat};
\addplot[mFR] table[x expr=\thisrow{h}*0.2, y=frdistens] {data/h_starlink_dl_rtt_MIS.dat};
\nextgroupplot[title={5G DL}]
\addplot[mRidge] table[x expr=\thisrow{h}*0.2, y=ridgefr] {data/h_quectel_dl_rtt_MIS.dat};
\addplot[mPrev] table[x expr=\thisrow{h}*0.2, y=prevseq2seqens] {data/h_quectel_dl_rtt_MIS.dat};
\addplot[mResid] table[x expr=\thisrow{h}*0.2, y=frdist] {data/h_quectel_dl_rtt_MIS.dat};
\addplot[mFR] table[x expr=\thisrow{h}*0.2, y=frdistens] {data/h_quectel_dl_rtt_MIS.dat};
\nextgroupplot[title={OneWeb UL}, ylabel={RTT interval score (ms)}]
\addplot[mRidge] table[x expr=\thisrow{h}*0.2, y=ridgefr] {data/h_oneweb_ul_rtt_MIS.dat};
\addplot[mPrev] table[x expr=\thisrow{h}*0.2, y=prevseq2seqens] {data/h_oneweb_ul_rtt_MIS.dat};
\addplot[mResid] table[x expr=\thisrow{h}*0.2, y=frdist] {data/h_oneweb_ul_rtt_MIS.dat};
\addplot[mFR] table[x expr=\thisrow{h}*0.2, y=frdistens] {data/h_oneweb_ul_rtt_MIS.dat};
\nextgroupplot[title={Starlink UL}]
\addplot[mRidge] table[x expr=\thisrow{h}*0.2, y=ridgefr] {data/h_starlink_ul_rtt_MIS.dat};
\addplot[mPrev] table[x expr=\thisrow{h}*0.2, y=prevseq2seqens] {data/h_starlink_ul_rtt_MIS.dat};
\addplot[mResid] table[x expr=\thisrow{h}*0.2, y=frdist] {data/h_starlink_ul_rtt_MIS.dat};
\addplot[mFR] table[x expr=\thisrow{h}*0.2, y=frdistens] {data/h_starlink_ul_rtt_MIS.dat};
\nextgroupplot[title={5G UL}]
\addplot[mRidge] table[x expr=\thisrow{h}*0.2, y=ridgefr] {data/h_quectel_ul_rtt_MIS.dat};
\addplot[mPrev] table[x expr=\thisrow{h}*0.2, y=prevseq2seqens] {data/h_quectel_ul_rtt_MIS.dat};
\addplot[mResid] table[x expr=\thisrow{h}*0.2, y=frdist] {data/h_quectel_ul_rtt_MIS.dat};
\addplot[mFR] table[x expr=\thisrow{h}*0.2, y=frdistens] {data/h_quectel_ul_rtt_MIS.dat};
\end{groupplot}
\node at ($(group c2r2.south)+(0,-1.25cm)$) {\pgfplotslegendfromname{legMIS}};
\end{tikzpicture}
}
  \caption{RTT 95\% interval score per horizon step, with intervals calibrated on the validation sessions. Lower is better.}
  \label{fig:mis}
\end{figure}

\begin{figure}[t]
  \centering
  \resizebox{\linewidth}{!}{\begin{tikzpicture}
\begin{axis}[frbase, width=0.95\linewidth, height=5.2cm, xlabel={time in test session (s)}, ylabel={Starlink DL RTT (ms)}, xmin=0, xmax=59.8, legend style={at={(0.5,1.03)}, anchor=south, legend columns=4}]
\addplot[name path=hi, draw=none, forget plot] table[x=t, y=frhi] {data/trace_starlink_dl.dat};
\addplot[name path=lo, draw=none, forget plot] table[x=t, y=frlo] {data/trace_starlink_dl.dat};
\addplot[cFR!25] fill between[of=hi and lo];
\addlegendentry{Final 95\% interval (CQR)}
\addplot[black, thin] table[x=t, y=truth] {data/trace_starlink_dl.dat};
\addlegendentry{measured}
\addplot[mPrev, mark=none] table[x=t, y=y2] {data/trace_starlink_dl.dat};
\addlegendentry{Initial (ens.)}
\addplot[cFR, thick] table[x=t, y=frmed] {data/trace_starlink_dl.dat};
\addlegendentry{Final median (ens.)}
\end{axis}
\end{tikzpicture}
}
  \caption{One minute of a Starlink DL test session, forecast 1\,s ahead ($h=5$). The final model's interval widens before and during the RTT bursts and narrows in quiet periods. A constant-width interval (Initial) has to be as wide as the bursts everywhere to reach the same coverage.}
  \label{fig:trace}
\end{figure}

\subsection{Comparison with the Modeling Dynamics report}
Table~\ref{tab:y2} puts the final results next to the numbers printed in the Modeling Dynamics report, using the earlier ``aggregate'' (the mean over the three targets and 15 steps). The retrained Initial model matches the MAE and RMSE of the Modeling Dynamics report closely. The $R^2$ differs on OneWeb UL, where the report averaged the two constant targets with $R^2=1$ into the aggregate (audit item A4), and on Starlink DL, where two of the three Initial seeds have a negative jitter $R^2$ and the causal ping inputs remove the leaked part of the signal (A3). With the initial interval method, which takes the quantiles of the test residuals, the initial and final models have almost the same interval score (within about 5\%). With intervals calibrated on held-out sessions, the final model is much better. The final advantage is therefore the input-dependent interval width, which the initial evaluation could not show.

\begin{table}[t]
\centering
\caption{initial style aggregate (mean over RTT, jitter, and loss and over the 15 steps). MIS (Initial method) uses quantiles of the test residuals, as in the Modeling Dynamics report. MIS (calibrated) uses intervals calibrated on validation sessions.}
\label{tab:y2}
\scriptsize
\begin{tabular}{llccccc}
\toprule
Link & Source & MAE & RMSE & $R^2$ & MIS (Initial method) & MIS (calibrated) \\
\midrule
\tabinput{data/tab_y2compare.tex}
\bottomrule
\end{tabular}
\end{table}

Table~\ref{tab:isproto} shows the same effect for RTT alone. With intervals taken from the test residuals around the mean, the two models are within about 2\% of each other on every link. With calibrated intervals the final model is \ChgMISWorst\% to \ChgMISBest\% better.

\begin{table}[t]
\centering
\caption{RTT interval score (ms, seed means) with the initial method (test-residual quantiles around the mean) and with intervals calibrated on validation sessions.}
\label{tab:isproto}
\small
\begin{tabular}{lrrrrrr}
\toprule
 & \multicolumn{3}{c}{Earlier method} & \multicolumn{3}{c}{Calibrated} \\
\cmidrule(lr){2-4}\cmidrule(lr){5-7}
Link & Initial & Final & $\Delta$\% & Initial & Final & $\Delta$\% \\
\midrule
\tabinput{data/tab_isproto.tex}
\bottomrule
\end{tabular}
\end{table}

\subsection{Which step helps?}
Table~\ref{tab:ablation} adds the final steps one at a time. The residual decoder alone (Step 1) and the causal pipeline with context features (Step 2) do not improve the MAE. On OneWeb DL they make it worse (by 9.8\% and 19.0\%), because a squared-error model anchored on a spiky last value follows the spikes. Step 2 reduces the RMSE slightly on most links (up to 4.5\% on Starlink DL), which shows that the context features carry information about the conditional mean. The gains in MAE and in the interval score come from Step 3, the distributional heads. The ensemble (Step 4) adds up to 6.4 further points (Starlink DL MAE) and less than one point on the 5G links.

\begin{table}[t]
\centering
\caption{Ablation: change (\%) of RTT MAE, RMSE, and calibrated interval score (MIS) relative to the retrained Initial model (seed means, with the last row comparing the Final ensemble with the Initial seed mean). Negative is better.}
\label{tab:ablation}
\scriptsize
\setlength{\tabcolsep}{2pt}
\resizebox{\linewidth}{!}{%
\begin{tabular}{l*{18}{r}}
\toprule
 & \multicolumn{3}{c}{OneWeb DL} & \multicolumn{3}{c}{OneWeb UL} & \multicolumn{3}{c}{Starlink DL} & \multicolumn{3}{c}{Starlink UL} & \multicolumn{3}{c}{5G DL} & \multicolumn{3}{c}{5G UL} \\
\cmidrule(lr){2-4}\cmidrule(lr){5-7}\cmidrule(lr){8-10}\cmidrule(lr){11-13}\cmidrule(lr){14-16}\cmidrule(lr){17-19}
Steps & MAE & RMSE & MIS & MAE & RMSE & MIS & MAE & RMSE & MIS & MAE & RMSE & MIS & MAE & RMSE & MIS & MAE & RMSE & MIS \\
\midrule
\tabinput{data/tab_ablation.tex}
\bottomrule
\end{tabular}}
\end{table}

\subsection{Simple references}
Table~\ref{tab:baselines} compares the neural models with the four references. Two facts were not visible in the initial study, which had no linear or naive baseline. First, persistence is a poor forecast for RTT at 5\,Hz: its RMSE is higher than that of the neural models and of Ridge-E on every link, because single pings are noisy. Second, a tuned linear model on the final features (Ridge-E) is close to the initial network on RMSE skill, within one point on four links. On Starlink DL a repeat of the value one second earlier already has 28\% skill, which reflects a one-second pattern in the downlink RTT under the 1\,Hz iperf load. The final ensemble has the lowest MAE on every link and the highest skill on every link, tied with Ridge-E and the single Final models on 5G UL. The plain ridge model on initial features fails on 5G UL ($-43\%$ skill) because of unclipped logger sentinels, which the final pipeline removes.

\begin{table}[t]
\centering
\caption{RTT MAE (ms) and RMSE skill against persistence (\%, higher is better) for the reference methods and the neural models (seed means).}
\label{tab:baselines}
\scriptsize
\setlength{\tabcolsep}{3pt}
\begin{tabular}{l*{12}{r}}
\toprule
 & \multicolumn{2}{c}{OneWeb DL} & \multicolumn{2}{c}{OneWeb UL} & \multicolumn{2}{c}{Starlink DL} & \multicolumn{2}{c}{Starlink UL} & \multicolumn{2}{c}{5G DL} & \multicolumn{2}{c}{5G UL} \\
\cmidrule(lr){2-3}\cmidrule(lr){4-5}\cmidrule(lr){6-7}\cmidrule(lr){8-9}\cmidrule(lr){10-11}\cmidrule(lr){12-13}
Method & MAE & skill & MAE & skill & MAE & skill & MAE & skill & MAE & skill & MAE & skill \\
\midrule
\tabinput{data/tab_baselines.tex}
\bottomrule
\end{tabular}
\end{table}

\subsection{Jitter and loss}
Table~\ref{tab:targets} gives the jitter and loss results. The point accuracy is similar for both models, with a few exceptions in both directions: Final has a higher loss MAE on OneWeb DL and Starlink UL (seed mean) and a lower one on Starlink DL and on both 5G links. The calibrated interval scores of Final are lower on every link and target, often by a large factor (for example Starlink DL jitter: 2.8 against 8.5). Jitter and loss are one-second iperf reports with long runs of zeros and rare large values, which is the case where input-dependent intervals help most.

\begin{table}[t]
\centering
\caption{Jitter (ms) and loss (\%) on the test sessions (seed means, OneWeb UL omitted because both are zero in every row).}
\label{tab:targets}
\scriptsize
\begin{tabular}{llcccccc}
\toprule
 & & \multicolumn{3}{c}{Jitter} & \multicolumn{3}{c}{Loss} \\
\cmidrule(lr){3-5}\cmidrule(lr){6-8}
Link & Model & MAE & RMSE & MIS & MAE & RMSE & MIS \\
\midrule
\tabinput{data/tab_targets.tex}
\bottomrule
\end{tabular}
\end{table}

\subsection{Effect of the interpolated Starlink targets}
About \InterpFracSLDL\% of the Starlink DL test targets are interpolated (audit item A3). Restricted to measured targets, the ensemble RTT MAE is \RealMAEFRSLDL\,ms for Final and \RealMAEPrevSLDL\,ms for Initial (\AllMAEFRSLDL\ and \AllMAEPrevSLDL\,ms on all targets). The conclusions do not depend on the interpolated targets.
